\documentclass{article}

\usepackage{amsmath}
\usepackage{enumitem}
\usepackage[boxed,vlined,boxruled]{algorithm2e}

\title{Bit Manipulation}
\author{Kanishk Goel}
\date{7 May, 2026}
\parindent 0px

\SetAlgoCaptionLayout{centerline}

\begin{document}
	\SetAlgoCaptionSeparator{}
	\renewcommand{\thealgocf}{} 
	\renewcommand{\AlCapNameFnt}{\bfseries}
	\renewcommand{\AlCapFnt}{}
	\DontPrintSemicolon
	\renewcommand{\algorithmcfname}{}
	\maketitle
	\SetAlgoNoEnd
	\section{Binary Numbers}
	Integers are represented as binary numbers by computers. Both signed and unsigned positive integers are represented with their binary digits and negative signed numbers through their Two's complement.
	\section{Tricks using bitwise operations}
	\begin{itemize}
		\item $1 \ll x$ is a number with only the $x^{th}$ bit set.
		\item $\sim (1 \ll x)$ is a number with only the $x^{th}$ bit not set.
		\item $n \, | (1 \ll x)$ sets the $x^{th}$ bit in the number $n$.
		\item $n \, \& \sim (1 \ll x)$ clears the $x^{th}$ bit in the number $n$.
		\item $n \land (1 \ll x)$ flips the $x^{th}$ bit in the number $n$.
		\item $(n \gg x) \, \& 1$ checks if the $x^{th}$ bit is set.
		\item $n$ is divisible by $2^k$ if $n \, \& \, (2^k - 1) = 0$
		\item a predecessor of a power of $2$ number $n$, is the number $n-1$. It can be used to check if an arbitrary number $y$ is a power of $2$, because bitwise and of the number with its predecessor will be $0$. So, $n$ is a power of $2$ if $n \, \& \, (n-1) = 0$
		\item A number $n$, can have its rightmost set bit cleared by $n \, \& \, (n-1)$.
		\item $n \, \& \, (n+1)$ clears all trailing ones $0011 \, 0111_2 \implies 0011 \, 0000_2$.
		\item $n \, | \, (n+1)$ sets the last cleared bit $0011 \, 0101_2 \implies 0011 \, 0111_2$.
		\item $n \, \& \, -n$ extracts the last set bit $ 0011 \, 0100_2 \implies 0000 \, 0100_2 $
	\end{itemize}
	\section{Brian Kernighan's algorithm}
	\textbf{Objective:} To count the number of set bits in the number $n$. 
	
	\vspace{1.5mm}
	\textbf{Idea:} Repeatedly clear the rightmost set bit until the number becomes $0$.
	
	\section{Counting set bits of all numbers upto n}
	\textbf{Idea: } For numbers upto $2^x$ i.e $(1 \, \text{to} \; 2^x - 1)$ there are $x \cdot 2^{x - 1}$	number of set bits. \newline
	
	This leads to the following algorithm:
	\begin{enumerate}
		\item Find the highest power of $2$ lesser than equal to $n$. Let it be $x$.
		\item Calculate the number of set bits using the above formula.
		\item Count the number of set bits in the MSB from $2^x$ to $n$ and add it.
		\item Subtract $2^x$ from $n$ and repeat the above process.
	\end{enumerate}
	
	\section{Bitmasks and Enumerating Submasks}
	A binary digit is used as a flag in bitmasking to denote the status or existence of a feature or trait. \newline
	
	Given a bitmask $m$, a submask $s$ is one whose set bits are a subset of the set bits of $m$. To enumerate through all submasks, the idea is simply to start $s = m$ and then modify $s$ as $(s-1) \, \& \, m$.  \newline 
	
	\textbf{Justification:} $s-1$ removes the rightmost set bit, and bitwise and with $m$, matches the set bits of $m$ and $s-1$. This algorithm, thereby yields descending order of submasks of $m$. \newline
	
	Biggest application of this is in bitmask dp. Manier times bitmask dp requires iteration through all bitmasks, and for each mask, iteration through all its submasks. This algorithm runs for $m = 0$ to $2^n$, for each $m$, the above submask algorithm is run. \newline 
	
	\textbf{Runtime:} $O(3^n)$ - If mask  $m$ has $k$ set bits, then it will have $2^k$ submasks. A total of $\binom{n}{k}$ such masks with $k$ set bits exist. The final expression becomes:-
	
	\[
	\sum^n_{k=0} \binom{n}{k} \cdot 2^k
	\]
	
	which by the binomial theorm equals to $3^n$.
\end{document}